\begin{pastexam}[label={px:ntust-110-ete-ds-04}]{臺科大電子所甲組 110 學年度 第 4 題}
（問答）

Suppose we have three functions:
\begin{minted}[fontsize=\small]{text}
foo(int N)
    do something in A steps
    for i equals 1 to N
        do something in B steps
        for j equals 1 to N
            do something in C steps

goo(int N)
    do something in A steps
    for i equals 1 to N
        do something in B step
        for j equals 1 to i
            do something in C steps

poo(int N)
    do something in A steps
    for i equals 1 to N
        do something in B steps
        for j equals 1 to 250
            do something in C steps
\end{minted}
\begin{enumerate}[label=(\alph*)]
\item (6\%) For each of the three methods above (foo(), goo(), and poo()), what is the exact
runtime for each method (in steps) in terms of A,B,C, and N? You should not count loop related
operations.
\item (4\%) For N = 100, which one is the fastest algorithm? Why?
\item (6\%) Assuming that A, B, and C are constants, what is the asymptotic runtime of each method
in terms of N?
\item (4\%) As N grows infinitely large, which one is the fastest algorithm?
\end{enumerate}
\end{pastexam}

\begin{pastsoln}
題目說不計迴圈本身的操作，所以只數三種 do something 各做幾次，再乘上各自的步數，和資結書
第~1.3.3~節的逐行計算一樣。 原題 goo 第 4 行的 ``B step'' 應該是 ``B steps''，不影響答案。

(a)
\begin{itemize}
\item foo：A 做 $1$ 次；外層 $i$ 從 $1$ 到 $N$，B 做 $N$ 次；內層每輪跑 $N$ 次，C 做
$N \cdot N$ 次。 $T_{\mathrm{foo}}(N) = A + NB + N^2 C$。
\item goo：內層 $j$ 從 $1$ 到 $i$，跑 $i$ 次，C 共做
$\sum_{i=1}^{N} i = N(N+1)/2$ 次。 $T_{\mathrm{goo}}(N) = A + NB + N(N+1)C/2$。
\item poo：內層固定跑 $250$ 次，C 共做 $250N$ 次。
$T_{\mathrm{poo}}(N) = A + NB + 250NC$。
\end{itemize}

(b) 代 $N = 100$：
\[
T_{\mathrm{foo}} = A + 100B + 10000C,\quad
T_{\mathrm{goo}} = A + 100B + 5050C,\quad
T_{\mathrm{poo}} = A + 100B + 25000C .
\]
三者的 A、B 部分相同，只比 C 的次數：$5050 < 10000 < 25000$，所以 goo 最快（假設
$C > 0$；$C = 0$ 時三者一樣）。

(c) A、B、C 是正的常數。 foo 的最高次項是 $CN^2$，是 $\Theta(N^2)$（資結書第~1.5~節：多項式只看
最高次項）。 goo 是 $\frac{C}{2}N^2 + (B + \frac{C}{2})N + A$，也是 $\Theta(N^2)$。 poo 是
$(B + 250C)N + A$，是 $\Theta(N)$。

(d) poo。 它是 $\Theta(N)$，另外兩個是 $\Theta(N^2)$，$N$ 夠大以後一定是 poo 最快。
交叉點也算得出來：poo 的 C 比 foo 少，等於 $250N < N^2$，也就是 $N > 250$；比 goo
少，等於 $250N < N(N+1)/2$，也就是 $N > 499$。 所以 $N \ge 500$ 時 poo 最快，
$N = 499$ 時 poo 和 goo 的 C 都做 $124750$ 次。 (b) 和 (d) 的答案不同，因為漸近的
比較只管 $N$ 夠大的情形，$N = 100$ 時常數 $250$ 還佔上風。
\end{pastsoln}
